% Part: lambda-calculus
% Chapter: syntax
% Section: beta

\documentclass[../../../include/open-logic-section]{subfiles}

\begin{document}

\olfileid[fr]{lam}{syn}{bet}
\olsection{\texorpdfstring{$\beta$}{β}-réduction}

Devant $(\lambd[m][(\lambd[y][y]) m])$,
il est naturel de supposer un lien avec
$\lambd[m][m]$ : le second terme semble
résulter d'une « simplification » du premier.
La notion de \emph{$\beta$-réduction}
formalise cette intuition.

Dans la suite, nous considérons les termes à
$\alpha$-équivalence près. La substitution sur
ces classes est bien définie par le
\olref[lam][syn][alp]{cor:sub}.

\begin{defn}[Contraction~$\beta$, $\bredone$]
\ollabel{defn:betacontr}
La \emph{contraction~$\beta$} ($\bredone$)
est la plus petite relation compatible sur
les termes qui satisfait la condition suivante :
\[
  (\lambd[x][N])Q \bredone \Subst{N}{Q}{x}.
\]
Nous disons que $P$ se \emph{contracte par~$\beta$}
en $Q$ si $P \bredone Q$. Un terme de la
forme $(\lambd[x][N])Q$ est appelé un
\emph{rédex}.
\end{defn}

\begin{prob} \ollabel{prob:def}
Donner les définitions inductives équivalentes
de la contraction~$\beta$, comme nous l'avons
fait pour le changement de variable liée dans la
\olref[lam][syn][alp]{defn:aconvone}.
\end{prob}

\begin{defn}[Réduction~$\beta$, $\bred$]
\ollabel{defn:betared}
La \emph{réduction~$\beta$} ($\bred$) est
la plus petite relation réflexive et transitive
sur les termes qui contient $\bredone$.
Nous disons que $P$ se \emph{réduit par~$\beta$}
en $Q$ si $P \bred Q$.
\end{defn}

Quand le contexte est clair, nous écrirons
$\redone$ pour $\bredone$, et $\red$
pour $\bred$. De manière informelle,
$M \bred N$ signifie que $M$ peut être
transformé en $N$ par zéro ou plusieurs
contractions~$\beta$.

\begin{defn}[$\beta$-normal]
Un terme qui ne contient plus aucun rédex
à contracter est dit \emph{$\beta$-normal}.
\end{defn}

Si $M \bred N$ et si $N$ est $\beta$-normal,
alors $N$ est une \emph{forme normale} de~$M$.
On peut se demander si cette forme est unique;
la réponse est oui, comme nous le verrons
plus loin.

Considérons quelques exemples.
\begin{enumerate}
\item Nous avons
\begin{align*}
(\lambd[x][xxy]) \lambd[z][z]
  & \redone (\lambd[z][z])(\lambd[z][z]) y \\
  & \redone (\lambd[z][z]) y \\
  & \redone y.
\end{align*}
\item « Simplifier » un terme peut en fait
  le rendre plus complexe :
\begin{align*}
(\lambd[x][xxy])(\lambd[x][xxy])
  & \redone (\lambd[x][xxy])(\lambd[x][xxy])y \\
  & \redone (\lambd[x][xxy])(\lambd[x][xxy])yy \\
  & \redone \dots
\end{align*}
\item Une contraction peut également laisser
  le terme inchangé :
\[
(\lambd[x][xx])(\lambd[x][xx])
\redone (\lambd[x][xx])(\lambd[x][xx]).
\]
\item Un terme peut aussi se réduire de
  plusieurs manières. Par exemple,
\[
(\lambd[x][(\lambd[y][yx]) z]) v
\redone (\lambd[y][yv]) z
\]
  en contractant l'application extérieure;
  tandis que
\[
(\lambd[x][(\lambd[y][yx]) z]) v
\redone (\lambd[x][zx]) v
\]
  résulte de la contraction intérieure.
  Dans ce cas, les deux termes obtenus se
  réduisent ensuite au même terme, $zv$.
\end{enumerate}

Le résultat commun du dernier exemple n'est
pas dû au hasard : il illustre une propriété
profonde du lambda-calcul, dite
\emph{propriété de Church--Rosser}.

\begin{digress}
En général, on peut réduire un terme de
plusieurs manières. On a donc proposé
différentes stratégies; la plus courante
est la \emph{stratégie naturelle}. Elle
contracte toujours le rédex le plus à gauche,
la position d'un rédex étant celle de son
premier symbole dans le terme. Elle possède
la propriété suivante : si une stratégie
quelconque conduit un terme à une forme
normale, la stratégie naturelle y conduit
également. Dans la suite, nous l'adopterons
sauf indication contraire.
\end{digress}

\begin{defn}[$\beta$-équivalence, $\equal$]
La \emph{$\beta$-équivalence} ($\equal$)
est la relation définie inductivement par
les règles suivantes :
\begin{enumerate}
\item $M \equal M$.
\item Si $M \equal N$, alors $N \equal M$.
\item Si $M \equal N$ et $N \equal O$,
  alors $M \equal O$.
\item Si $M \equal N$, alors $PM \equal PN$.
\item Si $M \equal N$, alors $MQ \equal NQ$.
\item Si $M \equal N$, alors
  $\lambd[x][M] \equal \lambd[x][N]$.
\item $(\lambd[x][N])Q \equal \Subst{N}{Q}{x}$.
\end{enumerate}
\end{defn}

Les trois premières règles rendent la
relation réflexive, symétrique et transitive;
les trois suivantes la rendent compatible;
la dernière y inclut la contraction~$\beta$.

De manière informelle, deux termes sont
$\beta$-équivalents si l'on peut passer de
l'un à l'autre par zéro ou plusieurs
contractions~$\beta$ ou par des étapes
inverses. Par définition, $M$ se contracte
en sens inverse par~$\beta$ vers $N$ si et
seulement si $N$ se contracte par~$\beta$
vers $M$.

Nous étendrons encore cette relation par
d'autres règles et noterons l'équivalence
ainsi obtenue $\equal[X]$, où $X$ est la
règle ajoutée.

\end{document}
